\chapter*{What Remains to Be Done}

The gaps in each part are cataloged in their respective sections.
Taken together, the most important open questions are these.

The \emph{density of replicating terms}
(Part~\ref{part:origins}, Conjecture~\ref{conj:density}) is
the lynch-pin of the origins of life argument.
If the replication machinery can always be added to a process
without disturbing its observable behavior, the phase transition
to life is inevitable.
If not, life requires special initial conditions.
This is the most urgent theorem to prove or refute.

The \emph{symplectic structure} of the bisimulation quotient
(Part~\ref{part:physics}) would make the analogy with classical
mechanics fully
rigorous.
Whether the space of bisimulation classes carries a natural
closed non-degenerate two-form --- whether the trace in the
reversibility construction is genuinely a momentum --- is the
deepest mathematical question raised by that part.

The \emph{universality of the reversibility envelope} --- whether
$\GSLT^\dagger$ is the \emph{minimal} reversible extension of
$\GSLT$ in the category of GSLTs --- would give the construction
a canonical status it currently lacks.

The \emph{colimit at limit ordinals} in the fibration must be
shown to exist and to inherit the physics of the stages below.
Without this, the transfinite tower is a formal object without
a limit.

And the \emph{hard problem} remains.
The framework offers structural analogues of sentience and
consciousness that are mathematically non-trivial and causally
efficacious.
Whether there is something it is like to be a learner reading its
own drift at a point in the lattice is a question the framework
does not answer.
We do not claim to have dissolved the hard problem.
We claim to have found the right coordinate system in which to
state it precisely.

\vspace{2em}
\noindent\hrule
\vspace{1em}

\noindent
The work is a beginning.
The six parts of the Turn, and the first half of the Prestige which
audits them, establish that a single mathematical structure --- the
graph-structured lambda theory, once it names a site of interaction and
meters it --- is rich enough to contain mind, physics, epistemology,
phenomenology, and biology as theorems rather than analogies.
Whether the structure is rich enough to contain \emph{everything}
observable is the question the framework was built to ask.
